\begin{bild}{Integrierte Präzisierung v1.6: ein Baukasten, mehrere Ausführungen}
Die Matrixbrücke in Kapitel~\ref{ch:v16-matrix} verbindet ganzzahlige Operationen und markierte $E_8$-Wurzeln. Zwei Gegenproben begrenzen ihre Reichweite: Ohne Ortsverschiebungen ergeben die inneren Walk-Matrizen die Identität (Kapitel~\ref{ch:v16-operations}); symmetrisch zulässige Energien können verschiedene Grundzustände wählen (Kapitel~\ref{ch:v16-state}). Die Unterscheidung zwischen Grammatik und Ausführung ist damit an konkrete Rechnungen gebunden.
\end{bild}
